%%%%%%%%%%%%%%%%
%%% gen 14 %%%
%%%%%%%%%%%%%%%%

\Chapter{14}

\begin{center}
\eProto{
THE INSCRIPTION OF ABRAM

THE HEBREW | MIGHTY KING

DESTROYER OF THE FIVE RE

BELLIOUS KINGS | BEFORE W

HOSE TERRIBLE SPLENDOR T

HOSE WHO HATE EL ELYON N

OW LIE SCATTERED AND SMAS

HED | THEIR LAND ALL IN ASH

ES | THEIR SEED IS NO MORE \cun{𒀭}
}
\end{center}

\Verse{1}
{
    \hOther{וַיְהִי בִּימֵי אַמְרָפֶל מֶלֶךְ שִׁנְעָר אַרְיוֹךְ מֶלֶךְ אֶלָּסָר
        כְּדרְלָעֹמֶר מֶלֶךְ עֵילָם וְתִדְעָל מֶלֶךְ גּוֹיִם׃}
}
{
    \eOther{And it was in the days of Amraphel, king of Shinar, Arioch, king of
        Ellasar, Chedorlaomer, king of Elam,
        and Tidal, king of Goiim.}
}
{
    \aC{
        \textsc{Regnal Year}: unspecified.

        \textsc{Theater}: Jordan Valley.

        \textsc{Eastern Coalition}: Amraphel of Shinar; Arioch of Ellasar; Chedorlaomer of Elam; Tidal of Goiim.
    }
}

\Verse{2}
{
    \hOther{עָשׂוּ מִלְחָמָה אֶת בֶּרַע מֶלֶךְ סְדֹם
        וְאֶת בִּרְשַׁע מֶלֶךְ עֲמֹרָה שִׁנְאָב מֶלֶךְ אַדְמָה
        וְשֶׁמְאֵבֶר מֶלֶךְ (צביים) [צְבוֹיִם]
        וּמֶלֶךְ בֶּלַע הִיא צֹעַר׃}
}
{
    \eOther{They made war with Bera, king of Sodom,
        and Birsha, king of Gomorrah, Shinab, king of Admah,
        and Shemeber, king of Zeboiim, and the king of Bela. (That is Zoar.)}
}
{
    \aC{\textsc{Opposing Coalition}: Bera of Sodom; Birsha of Gomorrah; Shinab of
        Admah; Shemeber of Zeboiim; king of Bela (modern designation: Zoar).
        Hostilities commenced between the two coalitions.}
}
{
    \aC{Qere Ketiv.}
}

\Verse{3}
{
    \hOther{כּל אֵלֶּה חָבְרוּ אֶל עֵמֶק הַשִּׂדִּים הוּא יָם הַמֶּלַח׃}
}
{
    \eOther{All these were allied at the Siddim Valley. (That is the Dead Sea.)}
}
{
    \aC{\textsc{Assembly Point}: Valley of Siddim (modern designation: Dead Sea). The
        five western kings concentrated their forces there.}
}
{

}

\Verse{4}
{
    \hOther{שְׁתֵּים עֶשְׂרֵה שָׁנָה עָבְדוּ אֶת כְּדרְלָעֹמֶר
        וּשְׁלֹשׁ עֶשְׂרֵה שָׁנָה מָרָדוּ׃}
}
{
    \eOther{Twelve years they served Chedorlaomer, and the thirteenth year they revolted,}
}
{
    \aC{\textsc{Political Status}: Twelve years under Chedorlaomer. \textsc{Year 13}: western vassals revolted.}
}
{

}

\Verse{5}
{
    \hOther{וּבְאַרְבַּע עֶשְׂרֵה שָׁנָה בָּא כְדרְלָעֹמֶר
        וְהַמְּלָכִים אֲשֶׁר אִתּוֹ וַיַּכּוּ אֶת רְפָאִים בְּעַשְׁתְּרֹת
        קַרְנַיִם וְאֶת הַזּוּזִים בְּהָם
        וְאֵת הָאֵימִים בְּשָׁוֵה קִרְיָתָיִם׃}
}
{
    \eOther{and in the fourteenth year Chedorlaomer
        and the kings who were with him came,
        and they struck the Rephaim in Ashteroth Karnaim
        and the Zuzum in Ham and the Emim in Shaveh Kiriataim}
}
{
    \aC{\textsc{Year 14}: Chedorlaomer and allied kings commenced punitive
        operations. Rephaim struck at Ashteroth Karnaim; Zuzim at Ham; Emim at
        Shaveh Kiriathaim;}
}
{

}

\Verse{6}
{
    \hOther{וְאֶת הַחֹרִי בְּהַרְרָם שֵׂעִיר עַד אֵיל פָּארָן אֲשֶׁר עַל
        הַמִּדְבָּר׃}
}
{
    \eOther{and the Horites in their mountain Seir to El Paran, which is by the
        wilderness.}
}
{
    \aC{Horites struck throughout the hill country of Seir as far as El
        Paran, on the wilderness frontier.}
}
{

}

\Verse{7}
{
    \hOther{וַיָּשֻׁבוּ וַיָּבֹאוּ אֶל עֵין מִשְׁפָּט הִוא קָדֵשׁ
        וַיַּכּוּ אֶת כּל שְׂדֵה הָעֲמָלֵקִי
        וְגַם אֶת הָאֱמֹרִי הַיֹּשֵׁב בְּחַצְצֹן תָּמָר׃}
}
{
    \eOther{And they came back and came to En Mishpat (that is Kadesh)
        and struck the area of the Amalekites
        and also the Amorites who live in Hazazon Tamar.}
}
{
    \aC{\textsc{Return Advance}: Coalition forces turned toward En Mishpat (modern
        designation: Kadesh), struck Amalekite territory, then the Amorites
        stationed at Hazazon Tamar.}
}
{

}

\Verse{8}
{
    \hOther{וַיֵּצֵא מֶלֶךְ סְדֹם וּמֶלֶךְ עֲמֹרָה
        וּמֶלֶךְ אַדְמָה וּמֶלֶךְ (צביים) [צְבוֹיִם]
        וּמֶלֶךְ בֶּלַע הִוא צֹעַר וַיַּעַרְכוּ אִתָּם מִלְחָמָה בְּעֵמֶק
        הַשִּׂדִּים׃}
}
{
    \eOther{And the king of Sodom and the king of Gomorrah
        and the king of Admah and the king of Zeboiim
        and the king of Bela (that is Zoar) went out
        and aligned with them for war in the Siddim Valley}
}
{
    \aC{\textsc{Western Response}: Kings of Sodom, Gomorrah, Admah, Zeboiim,
        and Bela (Zoar) deployed for battle in the Valley of Siddim.}
}
{
    \aC{Qere Ketiv.}
}

\Verse{9}
{
    \hOther{אֵת כְּדרְלָעֹמֶר מֶלֶךְ עֵילָם
        וְתִדְעָל מֶלֶךְ גּוֹיִם וְאַמְרָפֶל מֶלֶךְ שִׁנְעָר
        וְאַרְיוֹךְ מֶלֶךְ אֶלָּסָר אַרְבָּעָה מְלָכִים אֶת הַחֲמִשָּׁה׃}
}
{
    \eOther{with Chedorlaomer, king of Elam,
        and Tidal, king of Goiim, and Amraphel, king of Shinar,
        and Arioch, king of Ellasar---four kings with five.}
}
{
    \aC{\textsc{Order of Battle}: Chedorlaomer of Elam; Tidal of Goiim; Amraphel of
        Shinar; Arioch of Ellasar. \textsc{Strength}: four kings against five.}
}
{

}

\Verse{10}
{
    \hOther{וְעֵמֶק הַשִּׂדִּים בֶּאֱרֹת בֶּאֱרֹת חֵמָר
        וַיָּנֻסוּ מֶלֶךְ סְדֹם וַעֲמֹרָה
        וַיִּפְּלוּ שָׁמָּה וְהַנִּשְׁאָרִים הֶרָה נָּסוּ׃}
}
{
    \eOther{And Siddim Valley was pits, pits of bitumen,
        and the kings of Sodom and Gomorrah fled,
        and they fell there; and those who were left fled to the mountain.}
}
{
    \aC{\textsc{Terrain}: Valley of Siddim heavily broken by bitumen pits. \textsc{Result}:
        Sodomite and Gomorrite forces routed. Some fell among the pits;
        surviving elements withdrew into the hills.}
}
{

}

\Verse{11}
{
    \hOther{וַיִּקְחוּ אֶת כּל רְכֻשׁ סְדֹם
        וַעֲמֹרָה וְאֶת כּל אכְלָם וַיֵּלֵכוּ׃}
}
{
    \eOther{And they took all the property of Sodom
        and Gomorrah and all their food,
        and they went.}
}
{
    \aC{\textsc{Sodom and Gomorrah}: Coalition forces seized all movable property
        and provisions and withdrew.}
}
{

}

\Verse{12}
{
    \hOther{וַיִּקְחוּ אֶת לוֹט וְאֶת רְכֻשׁוֹ בֶּן אֲחִי אַבְרָם
        וַיֵּלֵכוּ וְהוּא יֹשֵׁב בִּסְדֹם׃}
}
{
    \eOther{And they took Lot, Abram's brother's son,
        and his property when they went.
        And he had been living in Sodom.}
}
{
    \aC{\textsc{Captives}: Lot, nephew of Abram
        and resident of Sodom, taken with his property during withdrawal.}
}
{

}

\Verse{13}
{
    \hOther{וַיָּבֹא הַפָּלִיט וַיַּגֵּד לְאַבְרָם הָעִבְרִי
        וְהוּא שֹׁכֵן בְּאֵלֹנֵי מַמְרֵא הָאֱמֹרִי אֲחִי אֶשְׁכֹּל
        וַאֲחִי עָנֵר וְהֵם בַּעֲלֵי בְרִית אַבְרָם׃}
}
{
    \eOther{And an escapee came and told Abram the Hebrew,
        and he was tenting among the oaks of Mamre the Amorite, brother of
        Eshcol and brother of Aner; and they were covenant partners of Abram.}
}
{
    \aC{
    	\textsc{Survivor Report}: One escapee reached Abram the Yid,%
		\footnote{
		    The term used here is not typically one Israelites use to describe
		    themselves, but is more characteristic of how they are described by or
		    among foreigners. This is one of several indications that the chapter
		    lies outside the primary sources of the Torah and may preserve or imitate
		    material belonging to the tradition of Mesopotamian royal inscriptions.
		}
    	then encamped by the oaks of Mamre the Amorite.
    	Mamre, Eshcol, and Aner were covenant allies of Abram.
    }
	%    \footnote{
	%        the Hebrew. This is an unusual use of the word “Hebrew.” Elsewhere in
	%        biblical stories it is used to identify Israelites only when one is
	%        speaking among foreigners. It is not the standard term for the people,
	%        which is rather “Israelite” at first and “Jew” later. Perhaps it is used
	%        here because there are not yet any other Israelites around, and Abraham
	%        himself is the foreigner. (Regarding the term “Hebrew slave,” see the
	%        comment on Exod 21:2.)
	%    }
}

\Verse{14}
{
    \hOther{וַיִּשְׁמַע אַבְרָם כִּי נִשְׁבָּה אָחִיו
        וַיָּרֶק אֶת חֲנִיכָיו יְלִידֵי בֵיתוֹ שְׁמֹנָה עָשָׂר
        וּשְׁלֹשׁ מֵאוֹת וַיִּרְדֹּף עַד דָּן׃}
}
{
    \eOther{And Abram heard that his brother had been taken prisoner,
        and he had his trained men, born in his house, unsheathe: three hundred
        eighteen. And he pursued as far as Dan.}
}
{
    \aC{\textsc{Mobilization}: Upon confirmation that his kinsman had been captured,
        Abram mustered 318 trained household-born retainers. \textsc{Pursuit}: north to
        Dan.}
}
{

}

\Verse{15}
{
    \hOther{וַיֵּחָלֵק עֲלֵיהֶם לַיְלָה הוּא
        וַעֲבָדָיו וַיַּכֵּם וַיִּרְדְּפֵם עַד חוֹבָה אֲשֶׁר מִשְּׂמֹאל
        לְדַמָּשֶׂק׃}
}
{
    \eOther{And he divided against them by night, he
        and his servants, and he struck them
        and pursued them as far as Hobah, which is at the left of Damascus.}
}
{
    \aC{\textsc{Night Action}: Abram divided his force, attacked with his retainers,
        routed the enemy, and pursued as far as Hobah, north of Damascus.}
}
{

}

\Verse{16}
{
    \hOther{וַיָּשֶׁב אֵת כּל הָרְכֻשׁ וְגַם אֶת לוֹט אָחִיו
        וּרְכֻשׁוֹ הֵשִׁיב וְגַם אֶת הַנָּשִׁים
        וְאֶת הָעָם׃}
}
{
    \eOther{And he brought back all the property,
        and he also brought back Lot, his brother,
        and his property, and also the women
        and the people.}
}
{
    \aC{\textsc{Recovery}: All seized property recovered. Lot
        and his property recovered. Women
        and other captives recovered.}
}
{

}

\Verse{17}
{
    \hOther{וַיֵּצֵא מֶלֶךְ סְדֹם לִקְרָאתוֹ אַחֲרֵי שׁוּבוֹ מֵהַכּוֹת אֶת
        כְּדרְלָעֹמֶר וְאֶת הַמְּלָכִים אֲשֶׁר אִתּוֹ אֶל עֵמֶק שָׁוֵה הוּא
        עֵמֶק הַמֶּלֶךְ׃}
}
{
    \eOther{And the king of Sodom came out to him after he came back from
        striking Chedorlaomer and the kings who were with him at the Shaveh
        Valley. (That is the valley of the king.)}
}
{
    \aC{\textsc{Return}: Following the defeat of Chedorlaomer
        and allied kings, Abram was met by the King of Sodom at the Valley of
        Shaveh (modern designation: King's Valley).}
}
{

}

\Verse{18}
{
    \hOther{וּמַלְכִּי צֶדֶק מֶלֶךְ שָׁלֵם הוֹצִיא לֶחֶם
        וָיָיִן וְהוּא כֹהֵן לְאֵל עֶלְיוֹן׃}
}
{
    \eOther{And Melchizedek, king of Salem, had brought out bread
        and wine. And he was a priest of El the Highest.}
}
{
    \aC{\textsc{Salem}: Melchizedek, King of Salem
        and priest of El Most High, brought out bread
        and wine.}
}
{
%    \footnote{
%        the Highest. Hebrew ‘elyn. This is how God is known at the end of the
%        Torah also (Deut 32:8). There (and in Ps 82:6) this is the epithet that
%        is used for God in connection with the formation of the nations of the
%        earth. That makes it notable that here the priest who is associated with
%        El the Highest is not from Abraham’s family but from another group.
%        Abraham is not pictured as the only worshiper of this God on earth.
%    }
}

\Verse{19}
{
    \hOther{וַיְבָרְכֵהוּ וַיֹּאמַר בָּרוּךְ אַבְרָם לְאֵל עֶלְיוֹן קֹנֵה
        שָׁמַיִם וָאָרֶץ׃}
}
{
    \eOther{And he blessed him and said, "Blessed is Abram to El the Highest,
        creator of skies and earth.}
}
{
    \aC{\textsc{Blessing}: "Blessed be Abram by El Most High, Maker of heaven
        and earth;}
}
{

}

\Verse{20}
{
    \hOther{וּבָרוּךְ אֵל עֶלְיוֹן אֲשֶׁר מִגֵּן צָרֶיךָ בְּיָדֶךָ
        וַיִּתֶּן לוֹ מַעֲשֵׂר מִכֹּל׃}
}
{
    \eOther{And blessed is El the Highest, who delivered your foes into your
        hand." And he gave him a tithe from everything.}
}
{
    \aC{and blessed be El Most High, who delivered your enemies into your
        hand." \textsc{Tribute}: Abram transferred one tenth of all recovered goods to
        Melchizedek.}
}
{

}

\Verse{21}
{
    \hOther{וַיֹּאמֶר מֶלֶךְ סְדֹם אֶל אַבְרָם תֶּן לִי הַנֶּפֶשׁ
        וְהָרְכֻשׁ קַח לָךְ׃}
}
{
    \eOther{And the king of Sodom said to Abram, "Give the persons to me,
        and take the property for you."}
}
{
    \aC{\textsc{Sodomite Claim}: King of Sodom requested return of the persons; Abram
        was authorized to retain the property.}
}
{

}

\Verse{22}
{
    \hOther{וַיֹּאמֶר אַבְרָם אֶל מֶלֶךְ סְדֹם הֲרִמֹתִי יָדִי אֶל}
    \hRed{יהוה}
    \hOther{אֵל עֶלְיוֹן קֹנֵה שָׁמַיִם וָאָרֶץ׃}
}
{
    \eOther{And Abram said to the king of Sodom, "I've lifted my hand to}
    \eRed{YHWH}
    \eOther{El the Highest, creator of skies and earth,}
}
{
    \aC{
    	\textsc{Abram's Oath}:
    	"I have raised my hand to
    }
    \eRed{YHWH}\eRed{\footnote{
		This verse is one of only 5 occurrences of the name YHWH outside
		the J source prior to the name being revealed in Exodus 3 and Exodus 6.

    	Overall, the word YHWH occurs about
    	3,190 times in the Primary History (Genesis through II Kings.)
    	1,829 times in the Torah (Genesis through Deuteronomy), and
		211 times from the beginning of Genesis to the end of Exodus 6.

		After Exodus 3, the E source begins using the name YHWH,
		while also continuing to use the term God (Elohim).

		After Exodus 6, the P source does the same.

		This first exception to the rule is not particularly interesting,
		given that the entirety of Genesis 14 appears not to have been
		written by any of the primary authors of the Torah.
		
		Regardless, the idea that two of the three primary authors of the
		Torah have this idea that the name YHWH was not known until
		the time of Moses in Egypt is a fact worth noting. This was
		among the earliest forms of evidence that eventually allowed
		us to determine how the Torah was written, and by whom.
    }}
    \aC{El Most High, creator of heaven and earth,}
}


\Verse{23}
{
    \hOther{אִם מִחוּט וְעַד שְׂרוֹךְ נַעַל}
    \hOther{וְאִם אֶקַּח מִכּל אֲשֶׁר לָךְ וְלֹא תֹאמַר אֲנִי הֶעֱשַׁרְתִּי אֶת אַבְרָם׃}

}
{
    \eOther{that, from a thread to a shoelace,
    		I won't take anything that is
        	yours, so you won't say,
        	'I made Abram rich.'}
}
{
    \aC{
	    I will take neither thread nor sandal strap
	    nor anything belonging to you, so that no
	    claim may be made that the King of Sodom
	    enriched Abram.
    }
}


\Verse{24}
{
    \hOther{בִּלְעָדַי רַק אֲשֶׁר אָכְלוּ הַנְּעָרִים
        וְחֵלֶק הָאֲנָשִׁים אֲשֶׁר הָלְכוּ אִתִּי עָנֵר אֶשְׁכֹּל
        וּמַמְרֵא הֵם יִקְחוּ חֶלְקָם׃}
}
{
    \eOther{Except only what the boys have eaten
        and the share of the people who went with me: Aner, Eshcol,
        and Mamre. They shall take their share."}
}
{
    \aC{\textsc{Exceptions}: provisions already consumed by Abram's men; contractual
        shares of allied commanders Aner, Eshcol,
        and Mamre. Their shares remain theirs."}
}
